\begin{lemma}
\label{lemma-limit-preserving-flat-at-primes}
In Situation \ref{situation-flattening-general}
assume $R \to S$ is essentially of finite presentation
and $M$ is an $S$-module of finite presentation. Let
$R' = \colim_{\lambda \in \Lambda} R_\lambda$
be a directed colimit of $R$-algebras. Let $I_\lambda \subset R_\lambda$
be ideals such that $I_\lambda R_\mu \subset I_\mu$ for all
$\mu \geq \lambda$ and set $I' = \colim_\lambda I_\lambda$.
If (\ref{equation-flat-at-primes}) holds for
$(R', I')$, then there exists a $\lambda \in \Lambda$ such that
(\ref{equation-flat-at-primes}) holds for $(R_\lambda, I_\lambda)$.
\end{lemma}

\begin{proof}
We first prove the lemma in case $R \to S$ is of finite presentation
and then we explain what needs to be changed in the general case.
We are going to write $S_\lambda = S \otimes_R R_\lambda$,
$S' = S \otimes_R R'$, $M_\lambda = M \otimes_R R_\lambda$, and
$M' = M \otimes_R R'$.
The base change $S'$ is of finite presentation over $R'$ and
$M'$ is of finite presentation over $S'$ and similarly for the
versions with subscript $\lambda$, see
Algebra, Lemma \ref{algebra-lemma-base-change-finiteness}.
By
Algebra, Theorem \ref{algebra-theorem-openness-flatness}
the set
$$
U' = \{\mathfrak q' \in \Spec(S') \mid
M'_{\mathfrak q'}\text{ is flat over }R'\}
$$
is open in $\Spec(S')$. Note that $V(I'S' + JS')$
is a quasi-compact space which is contained in $U'$ by assumption.
Hence there exist finitely many $g'_j \in S'$, $j = 1, \ldots, m$
such that $D(g'_j) \subset U'$ and such
that $V(I'S' + JS') \subset \bigcup D(g'_j)$.
Note that in particular $(M')_{g'_j}$ is a flat module over $R'$.

\medskip\noindent
We are going to pick increasingly large elements $\lambda \in \Lambda$.
First we pick it large enough so that we can find
$g_{j, \lambda} \in S_{\lambda}$ mapping to $g'_j$.
The inclusion $V(I'S' + JS') \subset \bigcup D(g'_j)$ means that
$I'S' + JS' + (g'_1, \ldots, g'_m) = S'$ which can be expressed as
$$
1 = \sum y_tk_t + \sum z_sh_s + \sum f_jg'_j
$$
for some $z_s \in I'$, $y_t \in J$, $k_t, h_s, f_j \in S'$.
After increasing $\lambda$ we may assume such an equation holds in
$S_\lambda$. Hence we may assume that
$V(I_\lambda S_\lambda + J S_\lambda) \subset \bigcup D(g_{j, \lambda})$. By
Algebra, Lemma \ref{algebra-lemma-flat-finite-presentation-limit-flat}
we see that for some sufficiently large $\lambda$ the modules
$(M_\lambda)_{g_{j, \lambda}}$ are flat over $R_\lambda$.
In particular the module $M_\lambda$ is flat over $R_\lambda$
at all the primes corresponding to points of
$V(I_\lambda S_\lambda + J S_\lambda)$.

\medskip\noindent
In the case that $S$ is essentially of finite presentation, we can write
$S = \Sigma^{-1}C$ where $R \to C$ is of finite presentation and
$\Sigma \subset C$ is a multiplicative subset. We can also write
$M = \Sigma^{-1}N$ for some finitely presented $C$-module $N$, see
Algebra, Lemma \ref{algebra-lemma-construct-fp-module}.
At this point we introduce $C_\lambda$, $C'$, $N_\lambda$, $N'$. Then in
the discussion above we obtain an open $U' \subset \Spec(C')$
over which $N'$ is flat over $R'$. The assumption that
(\ref{equation-flat-at-primes}) is true means that $V(I'S' + JS')$ maps
into $U'$, because for a prime $\mathfrak q' \subset S'$, corresponding
to a prime $\mathfrak r' \subset C'$ we have
$M'_{\mathfrak q'} = N'_{\mathfrak r'}$. Thus we can find
$g'_j \in C'$ such that $\bigcup D(g'_j)$ contains the image of
$V(I'S' + JS')$. The rest of the proof is exactly the same as before.
\end{proof}